\documentclass[a4paper,oneside,11pt]{book}

% Package dasar
\usepackage[utf8]{inputenc}     % untuk karakter UTF-8
\usepackage[T1]{fontenc}        % encoding font
\usepackage{newtxtext,newtxmath} % font Times New Roman
\usepackage{graphicx}
\usepackage{titling}
\usepackage{titlesec}
% \usepackage[backend=biber,style=ieee]{biblatex}
% \addbibresource{ref.bib}
% \setcounter{biburllcpenalty}{7000}
% \setcounter{biburlucpenalty}{8000}

% Package tambahan untuk layout dan formatting
\usepackage{geometry}           % untuk margin
\usepackage{setspace}           % untuk spasi baris
\usepackage{fancyhdr}           % untuk header/footer
\usepackage{hyperref}           % untuk link dan bookmark
\usepackage{xurl}               % agar URL panjang bisa wrap/pindah baris dengan rapi
\usepackage{listings}           % untuk code Python
\usepackage{xcolor}             % untuk warna
\usepackage{float}              % untuk posisi gambar
\usepackage{tocloft}            % Untuk membuat daftar isi menampilkan "Bab 1. Pendahuluan" instead of "1. Pendahuluan",
\usepackage{lipsum}             % untuk lorem ipsum
\usepackage{indentfirst}        % untuk indent paragraf pertama
\usepackage{enumitem}           % untuk kustomisasi penomoran list (A., B., dst)

% Pengaturan layout
\geometry{left=2cm, right=2cm, top=2cm, bottom=2cm}
\onehalfspacing  % spasi 1.5

% Penomoran section dengan huruf kapital (A, B, C, ...)
\renewcommand{\thesection}{\Alph{section}}
\renewcommand{\thesubsection}{\thesection.\arabic{subsection}}
\renewcommand{\thesubsubsection}{\thesubsection.\arabic{subsubsection}}

% Penomoran figure mengikuti section (misal, Gambar A.1, A.2, ...)
\renewcommand{\thefigure}{\thesection.\arabic{figure}}
\makeatletter
\@addtoreset{figure}{section}
\makeatother

\titleformat{\section}
  {\normalfont\fontsize{13pt}{15pt}\selectfont\bfseries}
  {\thesection}
  {0.5em}
  {}

\titlespacing*{\section}
  {0pt}
  {12pt}
  {3pt}

\titleformat{\subsection}
  {\normalfont\bfseries}
  {\thesubsection}
  {0.5em}
  {}

\titleformat{\subsubsection}
  {\normalfont}
  {\thesubsubsection}
  {0.5em}
  {}

% Format chapter
\titleformat{\chapter}[hang]
  {\normalfont\huge\bfseries\centering}{\chaptertitlename\ \thechapter.}{1em}{}

% Format code Python
% --- VS Code-like code style (dark) ---
\definecolor{vscodeBg}{HTML}{1E1E1E}
\definecolor{vscodeFg}{HTML}{D4D4D4}
\definecolor{vscodeKeyword}{HTML}{569CD6}
\definecolor{vscodeString}{HTML}{CE9178}
\definecolor{vscodeComment}{HTML}{6A9955}
\definecolor{vscodeNumber}{HTML}{858585}
\definecolor{vscodeRule}{HTML}{3C3C3C}

\lstdefinestyle{vscode-dark}{
  language=Python,
  basicstyle=\ttfamily\small\color{vscodeFg},
  backgroundcolor=\color{vscodeBg},
  keywordstyle=\color{vscodeKeyword}\bfseries,
  commentstyle=\color{vscodeComment},
  stringstyle=\color{vscodeString},
  numbers=left,
  numberstyle=\tiny\color{vscodeNumber},
  frame=single,
  rulecolor=\color{vscodeRule},
  breaklines=true,
  showstringspaces=false,
  columns=fullflexible,
  keepspaces=true,
  tabsize=2,
  xleftmargin=2cm,
  xrightmargin=1cm,
  captionpos=b,
  morekeywords={np, pd, plt, sns, array, DataFrame, Series, read_csv, read_excel, head, tail, info, describe, plot, scatter, bar, hist, fit, predict, transform, merge, concat, groupby, drop, fillna, dropna, iloc, loc}
}

\lstset{style=vscode-dark}

% Pengaturan hyperref
\hypersetup{
    colorlinks=true,
    linkcolor=black,
    filecolor=magenta,
    linktoc=all,
    urlcolor=blue,
    citecolor=red
}

% Info dokumen
\newcommand{\Pertemuan}{0}
\newcommand{\Materi}{MATERI}
\title{
  LAPORAN PRAKTIKUM\\
  PRAKTIKUM TEKNOLOGI INFORMASI DAN INSTALASI KOMPUTER\\
  PERTEMUAN \Pertemuan\\
  \Materi
}
\newcommand{\NamaMahasiswa}{John Doe}
\newcommand{\NIMMahasiswa}{11123456789}
\newcommand{\KelasMahasiswa}{X4}
\newcommand{\DosenPengampu}{Dr. Strange}

\author{\NamaMahasiswa\\\NIMMahasiswa}

\begin{document}

% Jika ingin daftar pustaka muncul meskipun belum ada entri di ref.bib,
% tambahkan entri BibTeX ke file ref.bib terlebih dulu.

% Title Page
\begin{titlingpage} 
  \begin{center}

    \vspace{4cm}
    \begin{Large} 
      \textbf{\thetitle} \\
    \end{Large}

    \vspace{3cm}
    
    \includegraphics[width=0.3\textwidth]{resource/logo-ugm.png}

    \vspace{3cm}
    
    Disusun oleh:

    \vspace{0.2cm}

    \begin{large}
      \begin{tabular}{p{4cm} c l}
          Nama           & : & \NamaMahasiswa  \\
          NIM            & : & \NIMMahasiswa   \\
          Kelas          & : & \KelasMahasiswa \\
          Dosen Pengampu & : & \DosenPengampu  \\
      \end{tabular}
    \end{large}
    
    \vfill

    \begin{Large}
      \textbf{PROGRAM STUDI D-IV TEKNOLOGI REKAYASA PERANGKAT LUNAK\\
      DEPARTEMEN TEKNIK ELEKTRO DAN INFORMATIKA\\
      SEKOLAH VOKASI\\
      UNIVERSITAS GADJAH MADA\\
      YOGYAKARTA\\
      2026}
    \end{Large}
    
  \end{center}
\end{titlingpage}

% Header/Footer
\pagestyle{fancy}
\fancyhf{}            % kosongkan header/footer default
\renewcommand{\headrulewidth}{0pt} % hilangkan garis header (atas)
\renewcommand{\footrulewidth}{0pt} % pastikan tidak ada garis footer
\fancyfoot[C]{\thepage} % taruh nomor halaman di footer tengah

% Pastikan halaman dengan style 'plain' (mis. ToC/List) juga tanpa garis atas
\fancypagestyle{plain}{
  \fancyhf{}
  \renewcommand{\headrulewidth}{0pt}
  \renewcommand{\footrulewidth}{0pt}
  \fancyfoot[C]{\thepage}
}


% --- Styling Judul Daftar Isi, Gambar, dan Kode ---
\renewcommand{\cfttoctitlefont}{\hfill\normalfont\Large\bfseries}
\renewcommand{\cftaftertoctitle}{\hfill}
\renewcommand{\cftbeforetoctitleskip}{1em}
\renewcommand{\cftaftertoctitleskip}{1.5em}

% --- Format Bab di Daftar Isi ---
\renewcommand{\cftchappresnum}{Bab~}
\renewcommand{\cftchapaftersnum}{.}
\renewcommand{\cftchapnumwidth}{3.5em}
\renewcommand{\cftchapfont}{\bfseries}
\renewcommand{\cftchappagefont}{\bfseries}

% --- Spasi dan Estetika ---
\setlength{\cftbeforechapskip}{0.8em} % jarak antar bab
\setlength{\cftbeforesecskip}{0.2em}  % jarak antar subbab
\renewcommand{\cftsecfont}{\normalfont}
\renewcommand{\cftsubsecfont}{\normalfont}

% --- Judul Lokal ---
\renewcommand{\contentsname}{\centering Daftar Isi}
\renewcommand{\listfigurename}{\centering Daftar Gambar}
\renewcommand{\lstlistlistingname}{\centering Daftar Kode}

% --- Gaya Daftar Gambar dan Daftar Kode ---
\renewcommand{\cftloftitlefont}{\hfill\normalfont\LARGE\bfseries}
\renewcommand{\cftafterloftitle}{\hfill}

\tableofcontents

\newpage
% Daftar Gambar
\listoffigures

\newpage
% Daftar Listing
\lstlistoflistings

% --- Halaman Laporan Pertemuan (seperti contoh gambar) ---
% Tidak pakai \chapter agar bisa menyatu dengan halaman sebelumnya.
\newpage
\begin{center}
  \textbf{
    \Large PERTEMUAN \Pertemuan\\
    {\Materi}
  }
\end{center}

\section{Tujuan Praktikum}
\begin{enumerate}
  \item Something
\end{enumerate}

\section{Dasar Teori}

\subsection{Something}
\par Lorem ipsum dolor sit amet, consectetur adipiscing elit. Sed do eiusmod tempor incididunt ut labore et dolore magna aliqua. Ut enim ad minim veniam, quis nostrud exercitation ullamco laboris nisi ut aliquip ex ea commodo consequat. Duis aute irure dolor in reprehenderit in voluptate velit esse cillum dolore eu fugiat nulla pariatur. Excepteur sint occaecat cupidatat non proident, sunt in culpa qui officia deserunt mollit anim id est laborum.


%contoh gambar dan kode
% \begin{figure}[H]
%     \centering
%     \includegraphics[width=0.5\textwidth]{resource/contoh-gambar.png}
%     \caption{Contoh Gambar}
%     \label{fig:contoh-gambar}
% \end{figure}

% \begin{lstlisting}[language=php, caption={Contoh Kode Python}, label={lst:contoh-kode}]
% def hello_world():
%     print("Hello, World!")
% \end{lstlisting}

\section{Hasil dan Pembahasan}

\par Sesuai dengan tujuan praktikum, berikut adalah hasil yang diperoleh dan pembahasan terkait materi yang telah dipelajari.

\section{Kesimpulan}

\par Link Google Colab: {\small\url{link-colab}}

\section{Daftar Pustaka}
% \printbibliography[heading=none]

\end{document}